% Scientific Reports research article.
% Compile: pdflatex main && bibtex main && pdflatex main && pdflatex main
% Supplementary Information is a separate document; the main text refers to it by explicit
% item number, and check_refs.py verifies every such reference before submission.
\documentclass[onecolumn,11pt]{article}
\usepackage[utf8]{inputenc}
\usepackage[T1]{fontenc}
\usepackage{amsmath,amssymb,amsfonts}
\usepackage{graphicx}
\usepackage{booktabs}
\usepackage[numbers]{natbib}
\usepackage[margin=1in]{geometry}
\usepackage{caption}
\usepackage[hidelinks]{hyperref}
\usepackage{float}
\usepackage{xcolor}
\usepackage{multirow}
\usepackage{lineno}
\linenumbers
\tolerance=1000
\emergencystretch=1em
% generated by assemble_supp.py - do not edit
\makeatletter
\newcommand{\supref}[1]{%
  \expandafter\ifx\csname suprefnum@#1\endcsname\relax
    \textbf{??SUPP:#1??}%
  \else\csname suprefnum@#1\endcsname\fi}
\newcommand{\suptab}[1]{Supplementary Table~S\supref{tab-#1}}
\newcommand{\supfig}[1]{Supplementary Fig.~S\supref{fig-#1}}
\expandafter\def\csname suprefnum@tab-cohorts\endcsname{1}
\expandafter\def\csname suprefnum@tab-perontology\endcsname{2}
\expandafter\def\csname suprefnum@tab-persistence\endcsname{3}
\expandafter\def\csname suprefnum@tab-modules_full\endcsname{4}
\expandafter\def\csname suprefnum@tab-highways\endcsname{5}
\expandafter\def\csname suprefnum@tab-highways_null\endcsname{6}
\expandafter\def\csname suprefnum@tab-highways_consecutive\endcsname{7}
\expandafter\def\csname suprefnum@tab-scgpt_highways\endcsname{8}
\expandafter\def\csname suprefnum@tab-cascades\endcsname{9}
\expandafter\def\csname suprefnum@tab-pmi_memory\endcsname{10}
\expandafter\def\csname suprefnum@tab-leiden_sweep\endcsname{11}
\expandafter\def\csname suprefnum@tab-hyperparam\endcsname{12}
\expandafter\def\csname suprefnum@tab-svd_sweep\endcsname{13}
\expandafter\def\csname suprefnum@tab-svd_capacity\endcsname{14}
\expandafter\def\csname suprefnum@tab-concepts\endcsname{15}
\expandafter\def\csname suprefnum@tab-seed_stability\endcsname{16}
\expandafter\def\csname suprefnum@tab-crossmodel\endcsname{17}
\expandafter\def\csname suprefnum@tab-evaluability\endcsname{18}
\expandafter\def\csname suprefnum@tab-ceiling\endcsname{19}
\expandafter\def\csname suprefnum@tab-cell_level\endcsname{20}
\expandafter\def\csname suprefnum@tab-sweep\endcsname{21}
\expandafter\def\csname suprefnum@tab-crossline\endcsname{22}
\expandafter\def\csname suprefnum@tab-causal_arms\endcsname{23}
\expandafter\def\csname suprefnum@tab-unannotated_new\endcsname{24}
\expandafter\def\csname suprefnum@tab-novel\endcsname{25}
\expandafter\def\csname suprefnum@tab-batch\endcsname{26}
\expandafter\def\csname suprefnum@fig-scgpt_highways\endcsname{1}
\expandafter\def\csname suprefnum@fig-umap_overview\endcsname{2}
\expandafter\def\csname suprefnum@fig-l11_detail\endcsname{3}
\expandafter\def\csname suprefnum@fig-annotation_projections\endcsname{4}
\makeatother

\bibliographystyle{unsrtnat}

\title{\textbf{Sparse autoencoders reveal organized biological knowledge but minimal
regulatory logic in single-cell foundation models}}
\author{Ihor Kendiukhov\textsuperscript{1,*}\\[6pt]
\textsuperscript{1}Institute of Medical Genetics and Applied Genomics, University of T\"ubingen, T\"ubingen, Germany\\
\textsuperscript{*}Correspondence: \texttt{ihor.kendiukhov@student.uni-tuebingen.de}}
\date{}

\begin{document}
\maketitle

\begin{abstract}
Whether single-cell foundation models encode causal regulatory logic, rather than statistical
co-expression, remains unclear. We trained TopK sparse
autoencoders on residual-stream activations from every layer of Geneformer V2-316M and scGPT
whole-human, producing atlases of 82,525 and 24,527 features, released as interactive platforms. The features are biologically organized: 29--59\% annotate to standard ontologies and form
co-activation modules. Under a capacity-matched comparison they are not hidden from linear decomposition---truncated SVD
needs rank 250--403 to match their reconstruction quality---but give a sparse parameterization
supporting per-feature intervention. Ablating one feature disrupts the genes it fires on by 35--67$\times$, but not the wider
process it is annotated with. Against CRISPRi perturbation in four cell lines,
only 1 of 66 transcription factors shows target-specific responses. Benchmarking against differential expression on the same cells, established network inference and
random controls explains why: that assay measures 6,546 genes and the median perturbed factor has
one curated target among them, so target-specific recovery is undefined for most of the panel. On perturbation data covering the transcriptome the
features do enrich for curated regulons, for 3 of 13 factors against 1 of 9 perturbations of
non-regulatory genes---the direction specificity predicts, on a panel too small to establish it
($p = 0.62$). Whether these representations carry a small amount of factor-specific
regulatory information is therefore not settled by currently available assays, and how much of a
regulon a screen can measure is essential to interpreting any claim about what a model encodes.
\end{abstract}

\section{Introduction}
Single-cell foundation models (scFMs) such as Geneformer~\cite{theodoris2023transfer} and
scGPT~\cite{cui2024scgpt} are used for cell-type annotation, perturbation-response prediction, and
gene-network inference. They are trained on millions of transcriptomic profiles and learn
contextual gene representations without explicit supervision on regulatory relationships. A central
question is whether these representations encode \emph{causal regulatory logic}---the directed
relationships between transcription factors (TFs) and their target genes---or only the
\emph{statistical co-expression} that correlates with, but does not constitute, regulation.

At the level of attention weights this question has been addressed by a companion
study~\cite{kendiukhov2025systematic}, which found that attention captures co-expression rather
than regulatory signal: trivial gene-level baselines outperform attention edges, pairwise scores
add zero incremental predictive value, and causal ablation of putatively regulatory heads produces
no behavioural effect. Attention weights, however, are one view among several. The residual
stream---the running sum of all layer outputs---may encode richer structure than attention patterns
alone reveal.

The superposition hypothesis~\cite{elhage2022superposition} offers a framework for this richer
structure. When a model must represent more concepts than it has dimensions, it encodes features as
nearly-orthogonal directions in activation space and activates only a sparse subset of them per
input. Sparse autoencoders
(SAEs)~\cite{olshausen1997sparse,sharkey2022taking,cunningham2023sparse,bricken2023monosemanticity}
resolve superposition by learning an overcomplete dictionary with a sparsity constraint,
decomposing dense activations into interpretable features. In large language models, SAE pipelines
have scaled to billions of features~\cite{gao2024scaling,templeton2024scaling}; they have not been
systematically applied to biological foundation models, where the ``concepts'' are genes, pathways,
and regulatory programs rather than linguistic entities.

Here we present a systematic SAE-based interpretability analysis of single-cell foundation models.
We train TopK SAEs~\cite{makhzani2013ksparse} on per-position residual-stream activations from all
18 layers of Geneformer V2-316M and all 12 layers of scGPT whole-human~\cite{cui2024scgpt},
producing atlases of 82,525 and 24,527 features respectively. The two models differ substantially
in architecture (Geneformer's rank-value tokenization vs.\ scGPT's dual gene-ID plus binned-value
embeddings, 18 vs.\ 12 layers, $d{=}1{,}152$ vs.\ $d{=}512$, bidirectional masked-token vs.\
specialized causal-attention generative training) and training data (30M vs.\ 33M cells), yet we
apply an identical analytical pipeline to both.

Two questions organize the analysis. The first is descriptive: what biological structure do these
dictionaries contain, how is it organized across layers, and how much of it is accessible to
standard linear decomposition of the same activations? The second is causal: do the features
encode which genes a transcription factor actually regulates? Answering the second question
quantitatively requires a reference point, because a low recovery rate is uninformative unless one
knows how much of the regulatory relationship is recoverable from the same measurements by any
method. We therefore benchmark the SAE against a perturbation-aware empirical ceiling---
differential expression of the CRISPRi knockdown itself---and against established
regulatory-inference baselines evaluated in an identical framework, and we replicate the
measurement in four cell lines.

To make these results accessible to the community, we release both atlases as interactive web
platforms---the Geneformer Feature Atlas
(\url{https://biodyn-ai.github.io/geneformer-atlas/}) and scGPT Feature Atlas
(\url{https://biodyn-ai.github.io/scgpt-atlas/})---enabling exploration of over 107,000 features
across 30 layers of two leading single-cell foundation models.

\section{Results}
\subsection*{An SAE feature atlas of Geneformer}
\label{sec:atlas}

We extracted per-position residual stream activations from all 18 layers of Geneformer V2-316M,
processing 2,000 non-targeting control cells from the Replogle genome-scale CRISPRi
dataset~\cite{replogle2022mapping} and obtaining 4,056,351 token positions per layer (mean 2,028
genes per cell; 336.4~GB in total). The control condition in this resource is defined by
perturbation label rather than by cell line, so the cohort spans all four lines it contains---591
Jurkat, 581 K562, 567 RPE1 and 261 HepG2 cells (\suptab{cohorts})---and the dictionaries
below are learned from that multi-line control population. For each layer we trained a TopK sparse
autoencoder with a
4$\times$ overcomplete dictionary (4,608 features from 1,152 input dimensions) and $k{=}32$
sparsity, subsampling 1~million positions per layer for training efficiency (Methods).

Table~\ref{tab:full18} presents training and annotation statistics for all 18 layers
(Fig.~\ref{fig:overview}). Variance explained peaks at layers 3--4 (85.2--85.3\%) and generally
declines thereafter to 76.8\% at layer~17, with a brief recovery at layers 10--11, indicating that
later-layer representations become progressively more distributed and harder to compress with a
fixed-size dictionary. Dead features---those never activating on 100K held-out positions---increase
from 0 at layer~0 to 70 at layer~16, with a trough at layers 9--11 (6--13 dead) suggesting a
mid-stack revival of structured representations. Dictionary directions are well separated
throughout: the mean absolute cosine similarity between decoder columns ranges from 0.033 to 0.040,
close to the theoretical floor for this many unit vectors in 1,152 dimensions (Methods).

\begin{table}[t]
\centering
\caption{\textbf{Geneformer SAE training and annotation statistics.} Each SAE uses a 4$\times$ overcomplete dictionary (4,608 features) with TopK $k{=}32$ sparsity. VarExpl = variance explained. Dead = $4{,}608 - \text{Alive}$. Aligned = features with $|\cos| > 0.5$ against any of the leading 50 principal axes; see \suptab{svd_capacity} for the capacity-matched comparison. MeanCos = mean absolute pairwise cosine between decoder vectors. Ann\% = fraction of alive features with $\geq$1 significant enrichment (FDR $< 0.05$). Enrichments = unique (feature, term) pairs. Per-database breakdown in \suptab{perontology}.}
\label{tab:full18}
\smallskip
\begin{tabular}{rccrcrccc}
\toprule
Layer & VarExpl & Alive & Dead & Aligned & Non-aligned & MeanCos & Ann\% & Enrichments \\
\midrule
0  & 83.9\% & 4,608 & 0   & 41 & 4,567 & 0.033 & 58.6\% & 23,959 \\
1  & 84.6\% & 4,606 & 2   & 27 & 4,579 & 0.033 & 57.4\% & 23,131 \\
2  & 84.8\% & 4,601 & 7   & 29 & 4,572 & 0.034 & 55.5\% & 23,383 \\
3  & 85.2\% & 4,595 & 13  & 12 & 4,583 & 0.035 & 56.4\% & 21,922 \\
4  & 85.3\% & 4,583 & 25  & 13 & 4,570 & 0.037 & 53.9\% & 19,910 \\
5  & 84.6\% & 4,576 & 32  & 8  & 4,568 & 0.036 & 52.1\% & 17,865 \\
6  & 83.3\% & 4,580 & 28  & 5  & 4,575 & 0.035 & 49.1\% & 16,716 \\
7  & 81.4\% & 4,591 & 17  & 6  & 4,585 & 0.036 & 47.7\% & 15,470 \\
8  & 80.4\% & 4,586 & 22  & 3  & 4,583 & 0.038 & 45.4\% & 15,679 \\
9  & 80.2\% & 4,595 & 13  & 4  & 4,591 & 0.036 & 50.3\% & 16,925 \\
10 & 81.0\% & 4,602 & 6   & 7  & 4,595 & 0.035 & 55.8\% & 19,587 \\
11 & 81.6\% & 4,598 & 10  & 4  & 4,594 & 0.037 & 56.2\% & 19,943 \\
12 & 81.0\% & 4,592 & 16  & 3  & 4,589 & 0.037 & 53.2\% & 19,105 \\
13 & 80.1\% & 4,583 & 25  & 3  & 4,580 & 0.038 & 50.7\% & 17,338 \\
14 & 80.0\% & 4,568 & 40  & 2  & 4,566 & 0.040 & 50.6\% & 17,719 \\
15 & 78.8\% & 4,543 & 65  & 5  & 4,538 & 0.038 & 49.0\% & 15,571 \\
16 & 76.9\% & 4,538 & 70  & 5  & 4,533 & 0.037 & 54.7\% & 16,080 \\
17 & 76.8\% & 4,580 & 28  & 12 & 4,568 & 0.039 & 47.0\% & 15,764 \\
\midrule
\textbf{Total} & & \textbf{82,525} & \textbf{419} & \textbf{189} & \textbf{82,336} & & & \textbf{336,067} \\
\bottomrule
\end{tabular}
\end{table}

The atlas comprises \textbf{82,525 alive features} across all 18 layers, of which \textbf{43,241}
(52.4\%) receive at least one significant ontology annotation against five databases (Gene Ontology
Biological Process~\cite{ashburner2000go}, KEGG~\cite{kanehisa2000kegg},
Reactome~\cite{jassal2020reactome}, STRING~\cite{szklarczyk2023string}, and
TRRUST~\cite{han2018trrust}).

\subsection*{An SAE feature atlas of scGPT}
\label{sec:scgpt_atlas}

To test whether the organization observed in Geneformer generalizes across architectures, we
applied an identical SAE pipeline to scGPT whole-human~\cite{cui2024scgpt}---a model trained on 33
million cells with fundamentally different design choices: a dual input of gene-identity tokens plus
binned expression-value embeddings (51 bins by default, in contrast to Geneformer's rank-value
tokenization), 12 transformer layers rather than 18, $d_\text{model}{=}512$ rather than 1,152, and a
generative causal-attention training objective rather than bidirectional masked-gene modelling.

We extracted per-position residual-stream activations from all 12 layers, processing 3,000 diverse
Tabula Sapiens cells (1,000 immune across 43 cell types, 1,000 kidney, 1,000 lung) and obtaining
3,561,832 token positions per layer. For each layer we trained a TopK SAE with a 4$\times$
overcomplete dictionary (2,048 features from 512 input dimensions) and $k{=}32$ sparsity.

Table~\ref{tab:scgpt_full12} presents training statistics for all 12 scGPT layers. Reconstruction
quality is higher than for Geneformer: variance explained ranges from 85.7\% (L7) to 93.5\% (L4),
mean 90.2\% against Geneformer's 81.7\%. Dead features are near zero: 49 across all 12 layers
(0.2\%), compared with 419 in Geneformer (0.5\%). The atlas comprises \textbf{24,527 alive
features}, of which \textbf{7,595} (31.0\%) are annotated against the same five databases. These
two atlases are built from different input distributions, so the comparison between them is
addressed separately under matched inputs below.

\begin{table}[t]
\centering
\caption{\textbf{scGPT SAE training and annotation statistics.} 4$\times$ dictionary (2,048 features), TopK $k{=}32$, 3.56M positions/layer. Notation as in Table~\ref{tab:full18}.}
\label{tab:scgpt_full12}
\smallskip
\begin{tabular}{rccrccc}
\toprule
Layer & VarExpl & Alive & Dead & MeanCos & Ann\% & Modules \\
\midrule
0  & 92.0\% & 2,027 & 21 & 0.038 & 32.7\% & 6 \\
1  & 93.0\% & 2,035 & 13 & 0.041 & 30.8\% & 6 \\
2  & 93.2\% & 2,038 & 10 & 0.040 & 28.7\% & 5 \\
3  & 93.2\% & 2,045 & 3  & 0.041 & 31.6\% & 7 \\
4  & 93.5\% & 2,048 & 0  & 0.041 & 30.4\% & 6 \\
5  & 92.1\% & 2,048 & 0  & 0.042 & 29.6\% & 7 \\
6  & 90.3\% & 2,048 & 0  & 0.046 & 28.9\% & 7 \\
7  & 85.7\% & 2,048 & 0  & 0.049 & 32.4\% & 6 \\
8  & 86.3\% & 2,048 & 0  & 0.046 & 28.9\% & 7 \\
9  & 86.9\% & 2,047 & 1  & 0.045 & 33.9\% & 5 \\
10 & 87.4\% & 2,047 & 1  & 0.044 & 31.7\% & 7 \\
11 & 88.6\% & 2,048 & 0  & 0.043 & 32.0\% & 7 \\
\midrule
\textbf{Total} & & \textbf{24,527} & \textbf{49} & & & \textbf{76} \\
\bottomrule
\end{tabular}
\end{table}

\subsection*{Biological annotation follows a U-shaped layer profile}
\label{sec:ushape}

Annotation of each feature's top-20 genes against five databases (Fisher's exact test,
Benjamini--Hochberg FDR $< 0.05$) reveals a U-shaped profile across layers
(Fig.~\ref{fig:ushape}). Annotation rates are highest at layers 0--1 (57--59\%), decline to a
minimum at layer~8 (45.4\%), recover to 55--56\% at layers 10--11, and decrease again at layers
16--17 (47--55\%). The scGPT atlas shows a weaker but qualitatively similar pattern across its 12
layers, from 28.7\% (L2) to 33.9\% (L9).

For Geneformer this pattern admits a functional reading in four zones.

\paragraph{Early layers (0--4): molecular machinery.} Features map cleanly onto existing ontology
terms, with the highest GO~BP (8,500--10,150 enrichments per layer), KEGG (2,050--2,650), and
Reactome (9,200--11,000) counts. STRING protein--protein interaction associations peak here
(248--302 per layer), as do TRRUST TF associations (133--164). Representative features encode
textbook biological programs (Table~\ref{tab:examples}).

\begin{table}[t]
\centering
\caption{\textbf{Representative SAE features.} Top genes = highest mean activation magnitude. Databases = ontology sources with significant enrichments.}
\label{tab:examples}
\smallskip
\small
\begin{tabular}{cp{3.8cm}p{3.5cm}l}
\toprule
Feature & Top Genes & Biological Identity & Databases \\
\midrule
\multicolumn{4}{l}{\textit{Layer 0 --- Molecular Programs}} \\
F3717 & CDK1, CDC20, DLGAP5, PBK & Cell cycle (G2/M) & GO, KEGG, React., STRING \\
F3607 & RRM1, E2F1, MCM4, MCM6, RAD51 & DNA replication / repair & GO, KEGG, React., STRING, TRRUST \\
F1116 & ARL6IP4, SQSTM1, JAK1 & B cell activation & GO, KEGG, Reactome \\
F4536 & DYNC1H1, TLN1, MYH9, SPTAN1 & Cytoskeleton / focal adhesion & GO, KEGG, Reactome \\
F2829 & TGFB1, PKN1, GADD45A, ZBTB7A & MAPK/TGF$\beta$ signaling & GO, KEGG, React., STRING, TRRUST \\
F1573 & MKI67, CENPF, TOP2A, CDK1, AURKB & Mitosis / chr.\ segregation & GO, KEGG, React., STRING \\
\midrule
\multicolumn{4}{l}{\textit{Layer 11 --- Integrative Programs}} \\
F2035 & \textit{(cell diff.\ genes)} & Cell diff.\ (neg.\ reg.) & GO, Reactome \\
F3692 & \textit{(ERAD pathway genes)} & ER-associated degradation & GO, KEGG, Reactome \\
F3933 & \textit{(signaling genes)} & Intracell.\ signaling (neg.\ reg.) & GO, Reactome \\
F2936 & \textit{(mitochondrial genes)} & Mitochondrion organization & GO, KEGG, Reactome \\
\bottomrule
\end{tabular}
\end{table}

\paragraph{Middle layers (5--9): abstract computation.} Annotation rates drop below 50\% at layers
6--8, consistent with intermediate representations that are harder to map onto a single ontology
term. GO~BP falls to 6,600--7,700 and STRING associations to 181--216 per layer.

\paragraph{Mid-late layers (10--12): re-specialization.} Annotation rates recover to 53--56\%, with
GO~BP returning to 8,200--8,800 and KEGG to 1,750--2,100. Features shift from molecular processes
to integrative cellular programs such as cell differentiation, intracellular signaling, and
organelle organization.

\paragraph{Terminal layers (15--17): prediction-focused.} A second annotation decline (47--55\%),
the largest number of dead features (28--70), and the most distributed representations. Features at
these layers respond broadly to perturbations but with less specificity than mid-layer features.

\subsection*{Features are rebuilt at every layer rather than carried forward}
\label{sec:crosslayer}

To understand how features relate across layers, we tracked decoder weight vector similarity from
layer~0 features to all subsequent layers. Features are overwhelmingly layer-specific: only 2--3\%
of features at any layer match a feature at the adjacent layer. From layer~0 the decay is steep:
114 matches at L1 (2.5\%), 67 at L4 (1.5\%), 10 at L8 (0.2\%), and effectively zero beyond L10
(\suptab{persistence}). No
layer~0 feature survives past layer~11.

Feature persistence anti-correlates with biological content. Of layer~0 features, 98.2\% are
transient (present in at most three layers), with 59.6\% annotated and a mean of 5.4 enrichment
terms. The rare moderate-persistence features (1.8\%, present in 4--10 layers) have only a 3.7\%
annotation rate and 0.1 mean enrichments. Biological content is carried by features rebuilt at
every layer, not by a persistent scaffold.

\subsection*{Features organize into co-activation modules}
\label{sec:modules}

To map the relational structure among features we constructed co-activation graphs at each of the
18 layers using pointwise mutual information (PMI~\cite{manning1999foundations}) over the 4~million
token positions, and identified communities with the Leiden algorithm~\cite{traag2019louvain}
(Methods).

Across all layers we identified \textbf{141 modules} (6--12 per layer) covering
\textbf{96.0--99.5\%} of alive features (\suptab{modules_full}). This is not a trivial outcome: with $k{=}32$ sparsity out
of 4,608 features, random co-activation would produce a connected but undifferentiated graph;
instead the partition yields well-separated communities with distinct biological identity.

Module count peaks at layer~5 (12 modules), suggesting maximal functional compartmentalization at
early-to-mid processing, and PMI edge density declines with depth (446K at L0 to 328K at L16),
paralleling the declining variance explained (Fig.~\ref{fig:modules}). At layer~0 six modules
correspond to cell cycle and DNA replication, immune signaling, metabolism, translation, protein
quality control, and cytoskeleton and adhesion. At layer~11 eight modules shift toward integrative
cellular programs: cell differentiation, intracellular signaling, mitochondrial organization,
vesicular transport, chromatin and transcription, protein modification, cell motility, and stress
response.

The scGPT atlas exhibits the same modular organization: 76 modules across 12 layers (5--7 per
layer), with a mean coverage of 96.3\%. Despite the smaller dictionary (2,048 against 4,608),
module count per layer is comparable (6.3 against 7.8 mean), suggesting that the number of
biological programs an SAE resolves scales sub-linearly with dictionary size.

Because a single Leiden resolution is a choice rather than a ground truth, we swept
$\gamma \in \{0.5, 0.75, 1.0, 1.5, 2.0\}$ at layers 0 and 11 (\suptab{leiden_sweep}); the graph
dimensions and memory footprint of the co-activation construction are in \suptab{pmi_memory}. The partition is not stable in a
strong partition-agreement sense---at $\gamma{=}0.5$ the graph collapses into two large
communities and at $\gamma{=}2.0$ it fragments into 48--56, with adjusted Rand indices of 0.17--0.26
between neighbouring resolutions. What is stable is module \emph{identity}: the same biological
themes are recognizable throughout the range, with larger $\gamma$ splitting one program into
finer subprograms. Module counts should therefore be read as the resolution-1.0 cut of a
hierarchical community structure.

\subsection*{Cross-layer information highways}
\label{sec:highways}

Given that features are almost entirely layer-specific, how does biological information flow
through the network? We computed cross-layer PMI between SAE feature activations, encoding the same
500,000 positions through both source- and target-layer SAEs. An information highway is a
source-layer feature with at least one dependency on a target-layer feature above a PMI threshold
$\tau$.

Under a TopK constraint two features can co-activate simply because of marginal activation rates,
so any fixed threshold on raw PMI risks inflating the highway count. We therefore computed a
permutation null for three long-range pairs (L0$\to$L5, L5$\to$L11, L11$\to$L17) by independently
column-shuffling the target-layer activation matrix 20 times, recomputing the statistic, and taking
the maximum across rounds of the 99th percentile of finite PMI as a per-pair null threshold. The
reported count uses $\tau = \max(3.0, \tau_{\text{null}}^{p99})$; the null threshold came out at
2.46 or below for every pair, so the effective threshold is 3.0 and the highway fractions are not
an artifact of the sparsity constraint.

For the long-range pairs (\suptab{highways}, \suptab{highways_null}), L0$\to$L5 is at 98.96\%,
L5$\to$L11 at 96.66\%, and L11$\to$L17 at 99.02\%, with mean maximum PMI between 6.61 and 6.79 and individual connections exceeding PMI~10.
Across all 17 consecutive pairs L$_{i}$$\to$L$_{i+1}$ the curve is near-flat in the 94.7--99.1\%
range (\suptab{highways_consecutive}) with no sharp drop at any transition (Fig.~\ref{fig:highways}); the slight taper at
L16$\to$L17 (94.70\%) is consistent with the terminal-layer specialization described above.
Information flow is therefore continuous across the Geneformer stack rather than bottlenecked at
particular layers---in line with layer-wise benchmarking of scFM representations in downstream
tasks~\cite{wang2025scdrugmap}.

Biologically meaningful cascades appear among the strongest cross-layer connections
(\suptab{cascades}). The mTORC1-regulation to autophagy connection (L0$\to$L5, PMI~=~9.55) recapitulates a well-known
signaling axis, and protein-modification features at L11 connect to angiogenesis regulation at L17
(PMI~=~10.62).

scGPT shows a different pattern (\suptab{scgpt_highways}, \supfig{scgpt_highways}): upstream
connectivity is consistently high (86.6--95.5\%), but
downstream connectivity falls from 95.7\% (L0$\to$L4) to 62.9\% (L8$\to$L11), indicating a
progressive concentration of information toward later layers that Geneformer does not show
(97.4--99.8\% downstream throughout).

\subsection*{Feature space geometry}
\label{sec:geometry}

Projecting decoder weight vectors with UMAP under cosine distance produces structureless point
clouds at every layer. This is the expected geometry rather than a failure of the projection:
decoder directions are quasi-orthogonal by construction, with a mean \emph{signed} pairwise cosine
of 0.0007 and a mean \emph{absolute} pairwise cosine of 0.033--0.040 (Table~\ref{tab:full18}), and
a within-module versus between-module effect size of Cohen's $d = 0.075$. Thousands of features
pack into 1,152 dimensions by spreading nearly uniformly over the hypersphere, so no
distance-preserving two-dimensional embedding of the directions themselves can be informative.

We therefore visualized the co-activation network directly using a force-directed layout
(\supfig{umap_overview}). Each module forms a spatially coherent community, with annotated
features distributed across all communities while sparsely annotated features concentrate centrally,
and the layer-11 graph resolves this structure in detail (\supfig{l11_detail}). Annotation-based projections (TF-IDF weighted ontology vectors) provide independent
confirmation that the module structure reflects biological similarity
(\supfig{annotation_projections}).

\subsection*{Cell type enrichment}
\label{sec:celltypes}

To connect features to cellular identity we computed, for each feature, its mean activation in
cells of each type among the 3,000 Tabula Sapiens cells spanning 56 cell types across immune,
kidney and lung tissue, and tested for enrichment (Fisher's exact test, Benjamini--Hochberg
correction).

In the scGPT atlas, 2,028 of 2,048 features (99.0\%) at layer~7 are enriched for at least one cell
type, with similar coverage at all layers. The patterns are tissue-coherent: immune features
activate preferentially in T cells, B cells, macrophages and dendritic cells; kidney features in
proximal tubule cells and podocytes; lung features in alveolar epithelial and pulmonary endothelial
cells. For the Geneformer atlas we passed the same Tabula Sapiens cells through the model and
encoded them with the multi-line control dictionaries; despite the mismatch in training
distribution---immortalized lines against primary tissue---the dictionaries generalize to these
tissue contexts. Features of both models tile cell-identity space
comprehensively, consistent with the models' training on diverse cell populations.

\subsection*{Batch and assay confounds}
\label{sec:batch}

Cell-type enrichment alone does not distinguish biology from donor or sequencing-assay batch
effects, a particular concern for Tabula Sapiens, which spans 24 donors and 4 sequencing
chemistries. For each scGPT feature at layers \{0, 4, 7, 11\} we computed the mutual information
between its activation and two categorical metadata fields: tissue (3 values) and cell type (56 values)
(\suptab{batch}), tissue serving as a coarse proxy for sample-level batch because the three tissues came
from different preparations. Across all four layers, mean mutual information with cell type is
about 2.5$\times$ that with tissue---at layer~11, 0.0040 against 0.0016---and essentially no
feature is tissue-dominant (0/1990 at L0, 0/2038 at L4, 0/1999 at L7, 1/1766 at L11). Coarse tissue
batch is therefore not the primary driver of scGPT feature activation. Donor and assay identity are
not recorded in the cached extraction metadata and a donor-level decomposition is beyond what these
data support.

\subsection*{Interactive feature atlases}
\label{sec:atlases}

To enable community access to the complete atlases we developed and deployed two interactive web
platforms: the Geneformer Feature Atlas
(\url{https://biodyn-ai.github.io/geneformer-atlas/}; 82,525 features across 18 layers) and the
scGPT Feature Atlas (\url{https://biodyn-ai.github.io/scgpt-atlas/}; 24,527 features across 12
layers). Both provide six integrated views---Layer Explorer, Feature Detail Pages, Module Explorer,
Cross-Layer Flow, Gene Search, and Ontology Search---are built with React, Vite and Plotly.js,
served via GitHub Pages, and require no installation.


\subsection*{What the dictionary adds over the principal subspace}
\label{sec:svd}

A dictionary of 4,608 atoms and a basis of 50 principal axes are not comparable objects, so the
observation that few features align with the leading axes cannot by itself establish that the
information is inaccessible to linear decomposition. We therefore computed the exact
eigendecomposition of the activation covariance at every layer and compared the two representations
under matched conditions (Fig.~\ref{fig:svd}, \suptab{svd_sweep} and \suptab{svd_capacity}).

\paragraph{Reconstruction-matched capacity.} Truncated SVD requires rank \textbf{250--403} to reach
the variance the SAE explains on held-out positions---403 at layer~0 (SAE 0.848), 317 at layer~5
(0.856), 308 at layer~11 (0.823) and 250 at layer~17 (0.775). A 50-axis basis is five- to eightfold
short of the SAE's own reconstruction capacity, so it is the wrong reference for asking what the
dictionary can see.

\paragraph{Where SAE directions sit.} For each unit-normalized decoder direction we computed the
cumulative projection $\rho_k$ onto the rank-$k$ principal subspace, for every $k$ up to the full
1,152 dimensions, against the random-direction expectation $\rho_k = k/d$. SAE directions are
systematically concentrated in the leading subspace: at $k{=}100$ the median direction has
$\rho_k = 0.21$ at layer~0 and 0.33 at layer~11, against 0.087 for a random direction, a two- to
fourfold excess. The median direction reaches half of its norm by $k = 247$ (L0), 190 (L5), 178
(L11) and 152 (L17), where a random direction requires 576, and reaches 90\% by $k = 664$--764
against 1,037. SAE features are not hidden from the principal subspace; they occupy its
high-variance part preferentially.

\paragraph{Low single-axis alignment is not a capacity artefact.} The fraction of features whose
absolute cosine with at least one of the leading $k$ axes exceeds 0.5 saturates by $k \approx
20$--50 and does not increase thereafter. Even against the \emph{complete} 1,152-dimensional basis
it is 0.98\% at layer~0, 0.13\% at layer~5, 0.09\% at layer~11 and 0.26\% at layer~17. Offering more
axes does not recruit more features, because each feature distributes its norm across many axes
rather than concentrating on one. The same conclusion holds across cosine thresholds
$\{0.3, 0.5, 0.7, 0.9\}$ (\suptab{svd_sweep}).

\paragraph{Annotation yield per direction.} We annotated the leading 50 principal axes with the
identical top-20-gene enrichment protocol used for SAE features, in both polarities, since a
principal axis is signed whereas a TopK feature is one-sided. Per direction, principal axes are at
least as annotatable as features: at layer~0, 98\% of axes carry at least one significant term with
a mean of 159.5 terms per direction and 2,029 unique terms, against 90\%, 64.4 and 1,977 for 100
randomly drawn SAE features; at layer~11 the corresponding numbers are 98\%, 178.8 and 2,083 against
93\%, 71.3 and 2,045. Biological annotation is therefore not the exclusive property of
non-aligned directions.

Taken together, these three measurements locate what the dictionary contributes. It is not access to
information that linear decomposition cannot reach---the leading principal subspace contains most of
each feature's norm and annotates readily. It is a different \emph{parameterization} of that
information: thousands of sparse, one-sided, individually interpretable directions in place of a few
hundred dense, signed, polysemantic ones. That is the property the rest of this study relies on,
because it is what makes per-feature annotation, per-feature co-activation structure, and
per-feature causal intervention possible at all.

\subsection*{How many distinct concepts the atlas contains}
\label{sec:concepts}

Dictionary atoms trained at different layers occupy different vector spaces, so their counts cannot
be pooled into a number of concepts held simultaneously in one representation. What can be compared
across layers is each feature's gene-level signature. We built a similarity graph over all 82,525
Geneformer features from their top-20 gene sets and partitioned it (Methods; Fig.~\ref{fig:concepts},
\suptab{concepts}).

The 82,525 atoms resolve into \textbf{18,711 distinct gene-level programs}, that is 4.4 atoms per
program. The count is insensitive to the clustering resolution: 18,695 at $\gamma{=}0.5$ and 18,748
at $\gamma{=}2.0$. The distribution is strongly bimodal: 17,690 programs are singletons---features
whose gene signature overlaps no other feature's---while 1,021 multi-atom programs absorb the
remaining 64,835 atoms, the largest containing 3,221. Redundancy is therefore concentrated in a
small number of heavily re-derived programs rather than spread evenly. As exact sets the signatures
are nearly all distinct: 82,442 distinct signatures, with only 56 shared by more than one feature.

Programs recur across depth. Of the multi-atom programs, 97.1\% appear at two or more layers, and
the mean program spans 3.16 layers. Read together with the near-complete representational turnover
between adjacent layers, this gives a consistent picture: each layer rebuilds a largely fresh
dictionary, but the gene-level programs those dictionaries express are re-derived throughout the
stack.

Within a single layer---the only place a features-per-space ratio is defined---the dictionary is
4$\times$ overcomplete (4,608 atoms in 1,152 dimensions, of which 4,598 are alive at layer~11) and
resolves 774--2,048 distinct programs depending on depth (1,301 at L0, 814 at L5, 1,252 at L11,
1,360 at L17). The packing is close to the geometric limit: mean absolute decoder coherence is
0.036 at layer~11 against a Welch bound of 0.0255 for 4,608 unit vectors in 1,152 dimensions, so
the dictionary is near-optimally spread rather than exploiting an unusually large capacity.


\subsection*{Causal ablation: what a feature is necessary for}
\label{sec:causal}

Annotation and co-activation establish that features correspond to biology, but not that the model
uses them. To test necessity we ablated single features at layer~11: the hidden state was encoded
through the SAE, one feature's activation set to zero, the state decoded and substituted, and the
forward pass continued to the output logits. Disruption is the mean absolute change in output
logits at a set of gene positions, and specificity is that quantity on a gene set divided by the
same quantity on all other positions (Fig.~\ref{fig:causal}, \suptab{causal_arms}).

The design of this test determines what it can show. A feature's annotation is assigned from the
genes it most strongly activates on, so evaluating disruption at those same genes asks whether
ablating a feature affects the genes that define it---a question whose answer is close to
tautological. We therefore partitioned the gene positions of each tested feature into three disjoint
sets: genes that are both in the feature's top-20 activating genes and in its annotated term; genes
of the annotated term that are \emph{not} among its top-20 (an evaluation set independent of the
annotation evidence); and all remaining positions. Each feature additionally received its own null:
a random gene set matched to its held-out set in size and in mean-activation decile. Features were
drawn from three arms---the most richly annotated features, a uniformly random sample of annotated
features, and a uniformly random sample of all alive features---so that the result cannot depend on
how features were selected. All 119 tested features come from K562 cells.

Ablation has a large and reproducible effect on the genes a feature activates on. Median
specificity on the top-20 genes against all other positions is 35.5$\times$ in the richly annotated
arm, 64.3$\times$ and 62.1$\times$ in the two random arms, with over 92\% of features exceeding
2$\times$ in every arm. Restricted to the genes that are both in the top-20 and in the annotated
term, the medians are 37.3$\times$, 56.9$\times$ and 66.9$\times$.

On the independent evaluation set the effect disappears. Median specificity on annotated-term genes
outside the feature's top-20 is \textbf{1.04}, \textbf{0.94} and \textbf{0.97} in the three arms.
The matched random gene sets give 0.89, 0.91 and 0.89, so the held-out sets are statistically
indistinguishable from gene sets chosen at random with the same size and activation level: a
per-feature Wilcoxon signed-rank test of held-out against matched random gives $p = 0.76$ and
$p = 0.34$ in the two unbiased arms, and $p = 0.006$ with a median difference of 0.13 in the richly
annotated arm. Fourteen of the 119 features admit no independent evaluation at all, because every
gene of their annotated term already lies within their own top-20.

The conclusion is narrower than the ablation magnitudes alone suggest. Single-feature ablation
produces a specific, reproducible, gene-level effect, and that effect is confined to the genes on
which the feature fires. It does not extend to the remainder of the biological process the feature
is annotated with. Feature-level interventions therefore reveal genuine computational structure
that component-level ablation misses~\cite{kendiukhov2025systematic}, but the structure they reveal
is the feature's own support rather than the pathway it is labelled with.


\subsection*{Cross-model comparison on matched inputs}
\label{sec:crossmodel}

The two atlases were built from different cells, so any direct comparison of their statistics
confounds architecture with input distribution. To separate the two we trained Geneformer
dictionaries on the same 3,000 Tabula Sapiens cells used for the scGPT atlas, with identical
architecture and hyperparameters, and compared the models at matched relative depth
(\suptab{crossmodel}). Because the scGPT dictionaries were originally trained on all available
positions rather than a subsample, we also retrained them under the Geneformer protocol as a
control.

On matched inputs there is no uniform architectural advantage. Geneformer reconstructs better at
the input end (0.7731 against 0.7258 at relative depth 0), the two are indistinguishable at about a
quarter depth (0.7749 against 0.7729), and scGPT is ahead through the second half, reaching 0.8862
against 0.7241 at the output end. The difference between the models is a depth-dependent crossover
rather than a constant offset, and it is invisible unless the inputs are matched.

Input distribution accounts for a difference of the same order. Each Geneformer dictionary
evaluated on the distribution it was trained on reconstructs the CRISPRi-derived control cohort
4.7--7.8 percentage points better than the Tabula Sapiens dictionary reconstructs Tabula
Sapiens---0.8444 against 0.7731 at layer~0, for example---so a difference of this size between two
models is what differing inputs alone can produce. Both dictionaries sit close to the packing
limit for their dimensionality: mean absolute decoder coherence is 0.030--0.036 for Geneformer and
0.039--0.047 for scGPT, against Welch bounds of 0.0255 and 0.0383 respectively.

\subsection*{Robustness to training choices}
\label{sec:robustness}

\paragraph{Architecture.} Grid-training at layer~11 over expansion ratio $\{2\times, 4\times,
8\times\}$ and sparsity $k \in \{16, 32, 64\}$ gives variance explained rising monotonically with
both, from 0.706 to 0.884, with the reported configuration in the middle at 0.819
(\suptab{hyperparam}). Module count stays between 6 and 12 across the entire grid and mean decoder
coherence between 0.033 and 0.042, so the reported structure is not a property of one corner of the
hyperparameter space.

\paragraph{Random seed.} The atlas dictionaries were trained with seed 42. Retraining at layers 0
and 11 under six runs each---five differing only in initialisation and batch order, one with a
different training subsample---separates what is a property of the model from what is a property of
the run (\suptab{seed_stability}).

Everything the atlas reports in aggregate is reproducible. Variance explained is
$0.8454 \pm 0.0003$ at layer~0 and $0.8200 \pm 0.0003$ at layer~11, the annotation rate is
$0.935 \pm 0.009$ and $0.933 \pm 0.010$, module counts are $7.2 \pm 0.4$ and $8.3 \pm 1.0$, and
dead-feature counts are $2 \pm 1$ and $32 \pm 4$.

Individual dictionary atoms are not. The mean best decoder cosine between runs is 0.282 at layer~0
and 0.294 at layer~11, and fewer than 1\% of atoms have a counterpart above 0.9 in another run. The
reference for these numbers is what a random direction achieves: the best match a random unit
direction finds among 4,608 directions in 1,152 dimensions is $0.113 \pm 0.008$, and no random
direction reaches 0.9. Cross-seed agreement is therefore well above chance and far below
reproducibility. Independent runs learn different bases that explain the same variance and carry the
same annotation. This bears on how the released atlases should be read: their aggregate structure
is a property of the model, while the identity of any individual catalogued feature is partly a
property of the training run---consistent with the gene-level analysis above, in which the recurring
unit across layers is the program rather than the atom.

\subsection*{Features without ontology annotation}
\label{sec:unannotated}

Between 41\% and 55\% of features carry no significant ontology annotation, and it matters whether
these are structure the databases do not cover or noise. Clustering unannotated features by
top-20 gene-set similarity places only 2--3\% into standalone gene-set clusters, among them
ribosomal protein programs and mitochondrial complex assembly (\suptab{novel}).

Two tests separate them from annotated features (\suptab{unannotated_new}). Cell-type specificity
does not: at layers $\{0, 5, 11, 17\}$, 89.4--95.0\% of unannotated features are enriched for at
least one cell type against 94.5--95.8\% of annotated ones, and a matched random-feature null sits
at 91.9--94.7\%, so essentially every feature at these layers is cell-type selective and the test
has no discriminating power. Perturbation responsiveness does: at layer~11, 2.32\% of unannotated
features appear in a perturbation response list against 0.15\% of annotated features, a fifteenfold
difference, with unannotated features more responsive at layers 0 and 5 and indistinguishable at layer 17.

A cautious reading is that some unannotated features carry perturbation-response signal the
ontology databases do not describe. Being non-random is not the same as being biologically
meaningful, however: a feature that is cell-type selective may equally be tracking donor or assay
structure, which is why the batch analysis above is reported alongside.


\subsection*{What the perturbation data can support}
\label{sec:evaluability}

The preceding sections establish that features are biologically annotated, modularly organized and
causally specific. The question that remains is whether they encode regulatory logic: when a
transcription factor is knocked down, do the features that respond correspond to that factor's
known targets?

A recovery rate answers that question only against a reference. Before measuring the atlas we
therefore asked how much curated regulatory structure the assay can express at all. The CRISPRi
resource measures 6,546 genes. Of the 73 perturbed transcription factors with at least 50 cells,
only nine have five or more TRRUST targets among the measured genes, three have ten or more, and
one has twenty or more; the median perturbed factor has exactly \textbf{one} measurable target
(Fig.~\ref{fig:ceiling}b, \suptab{evaluability}). DoRothEA, whose regulons are larger, does
better---13 factors with five or more measured targets at the A+B+C confidence tier, median 15---but
the constraint is severe under every reference. For most of the panel, target-specific recovery is
not merely difficult; it is undefined, for any method.

\subsection*{A recoverability benchmark}
\label{sec:ceiling}

We therefore evaluated the SAE alongside methods that span the range of what is possible on these
data, each emitting a ranked list of 100 putative targets per factor and each scored by the same
hypergeometric enrichment with Benjamini--Hochberg correction across the panel and the measured
genes as the universe (Fig.~\ref{fig:ceiling}a, \suptab{ceiling}).

Differential expression of the knockdown itself is the perturbation-aware reference: it uses the
same cells and the actual intervention, and therefore bounds what any observational method can
achieve here. It recovers curated targets for 1/9 factors against TRRUST (11.1\%), 3/13 against
DoRothEA A+B+C (23.1\%) and 8/18 against DoRothEA at all confidence levels (44.4\%). GENIE3 reaches
11.1\%, 7.7\% and 0.0\% under the three references, GRNBoost2 11.1\%, 7.7\% and 5.6\%, Pearson
co-expression 22.2\%, 15.4\% and 11.1\%, and Spearman 22.2\%, 15.4\% and 5.6\%. Contextual gene
embeddings taken directly from the model, without any sparse decomposition, reach 11.1\%, 7.7\% and
0.0\%, and the SAE feature response recovers targets for no factor under any reference (0/9, 0/13,
0/18). Size-matched random gene sets recover nothing either, confirming that the test is
calibrated.

Two things follow. On the best-powered reference no observational method approaches the
intervention: 44.4\% against 0--11.1\%. Under the sparser references the ceiling is itself so
low---11.1\% under TRRUST---that co-expression matches or exceeds it, which is what a median of one
measurable target per factor produces: the reference is too thin to rank methods at all.
The choice of universe matters for this calibration: scoring against a nominal 20,000-gene universe
rather than the 6,546 measured genes makes random gene sets appear significant for 16.7\% of
factors, because predictions can only be drawn from what is measured.

\subsection*{Perturbation responses are detectable but not transcription-factor specific}
\label{sec:perturbation}

We mapped perturbation responses at layer~11 using up to 200 CRISPRi cells per target against 600
non-targeting cells of the same line, encoding every cell through the atlas dictionary (Methods).
The panel comprises 20 perturbed transcription factors, 18 of them evaluable, and as a
negative control ten
perturbations of genes that are not transcription factors under any reference.

\paragraph{The unit of replication.} A perturbation is applied per cell, but a per-position
statistic treats each gene position as an observation. The intraclass correlation of feature
activation across positions within a cell is 0.000298, which at the observed 2,020 positions per
cell gives a design effect of 1.61: the effective sample size is about 62\% of the nominal position
count. The correction is worth making and we make it throughout---significance is assessed over
cells---but its magnitude is bounded, because TopK sparsity leaves activations at different
positions of one cell close to independent.

\paragraph{Responses are weak and non-specific.} Thresholding on effect size alone---a shift
exceeding half a control standard deviation, with no correction for testing 4,608 features---the
panel yields a mean of 3.20 responding features per factor, with 15 of 20 factors showing at least
one. Correcting for multiple testing across features collapses this to two of twenty, so most of
what a per-feature threshold registers does not survive the number of features being tested. Under
that correction the like-for-like comparison of the two units of replication gives 1.60 responding
features per factor pooling positions and 1.90 assessing over cells (\suptab{cell_level}). The ten non-transcription-factor controls yield
\emph{more} responding features than the transcription factors (2.30 versus 1.90), so what the
features register is not specific to transcription-factor perturbation.

\paragraph{Where the sensitivity ends.} The clearest measurement of the limit compares the two
readouts on identical cells. Differential expression detects the knocked-down gene itself in 15 of
16 targets, at FDR $< 0.05$ with a median Cohen's $d$ of $-1.19$ (range $-2.76$ to $-0.49$). The
SAE feature response detects it in 1 of 20. The perturbation is therefore unambiguously present in
the expression profile of the same cells and is not registered by the feature responses. The
features are not insensitive because the data are thin; they track global cell-state shifts rather
than the gene-level consequences of the intervention. This is the specific sense in which detection
is limited, and it is a property of what the dictionary encodes.

\paragraph{Sample size.} Restricting to the 14 factors with at least 100 cells, so that sample size
is not confounded with which factors are available, the responding-feature count is flat across a
tenfold increase in cells (0.56 at $n{=}10$, 0.65 at $n{=}20$, 0.65 at $n{=}50$, 0.61 at
$n{=}100$). Detection of the knockdown rises from 5.0\% to 7.1\% and saturates. Nominal target
enrichment improves from $n{=}10$ to $n{=}50$ (7.7\% to 27.5\%) and then plateaus, the gain
saturates well before the effect becomes substantial, and under correction across the panel no
sample size yields a significant factor
(\suptab{sweep}).

\paragraph{Widening the training distribution.} A dictionary trained on the same control cells
pooled with 3,000 Tabula Sapiens cells from three primary tissues gives 2.55 responding features per
factor against 1.90 for the atlas dictionary, with the same non-transcription-factor controls at
2.30 and detection of the knockdown at 0 of 20. Diversifying the training distribution beyond
immortalized lines does not make regulatory structure recoverable.

\subsection*{Replication across cell lines and independent datasets}
\label{sec:replication}

The measurement was repeated in three further cell lines from the same resource, with perturbed and
non-targeting cells drawn from the same line in every case, and on two perturbation datasets
generated by other groups with other protocols and substantially larger gene panels
(Fig.~\ref{fig:power}c, \suptab{crossline}).

Across K562, RPE1, Jurkat and HepG2 the pattern is the same. Pooling the four lines,
\textbf{1 of 66} evaluable transcription factors shows target enrichment at FDR $< 0.05$
(1.5\%): none of 18 in K562, none of 16 in RPE1, one of 16 in Jurkat and none of 16 in HepG2. The
knocked-down gene itself is detected in 2 of 68 factors across the four lines. Responding-feature
counts are low throughout (means of 0.06--1.90 per factor) and do not separate transcription
factors from the non-transcription-factor controls, which average 0.17--2.30 in the same lines. The
design effect is 1.61--1.73 everywhere, so the correction for the unit of replication behaves
consistently across backgrounds. RPE1 is a non-transformed retinal pigment epithelial line and
HepG2 a hepatocyte line, so the result is not confined to the leukaemic backgrounds in which these
models are most often probed.

Both external datasets measure far more of the transcriptome than the primary resource, and the
difference changes the result. In THP-1 monocytes 18,649 genes are measured and the median
perturbed factor has 19.5 curated TRRUST targets among them; in the K562 Perturb-seq data 33,694
genes are measured. The Perturb-seq screen is a CRISPR \emph{activation} experiment, so it also
tests the question under gene induction rather than repression: the median perturbed factor is
induced with a Cohen's $d$ of $+0.98$ and yields 405 differentially expressed genes.

On these data the features do recover targets. Against DoRothEA the SAE reaches 20--22\% of factors
in the THP-1 data and 23--25\% in the Perturb-seq data, where on the primary panel it reached none.
Against TRRUST it reaches none in either, while differential expression on the same cells reaches
18\% and 30\%.

Whether that recovery is transcription-factor-specific requires a control, and we ran one: ten
perturbations of genes absent from every reference network, put through the identical pipeline,
with their predicted sets tested against the \emph{same} regulons as the transcription factors.

Against the high-confidence DoRothEA subset, a factor's own predicted set is enriched for its own
regulon in 3 of 13 cases (23\%). Of the nine control perturbations that produced any prediction,
one is enriched for at least one factor's regulon (11\%), giving 1 significant test in 117
factor-by-regulon comparisons. The difference is in the direction specificity would predict, but a
13-factor panel cannot establish it: Fisher's exact test on 3 of 13 against 1 of 9 gives
$p = 0.62$.

Against the union of all confidence levels the comparison is decisive in the other direction. The
matching factors reach 4 of 16 (25\%) while 8 of the 9 control perturbations enrich for at least
one regulon ($p = 0.004$, controls higher). Those regulons have a median of 402 measured genes and
reach 9,321, and a broad expression signature enriches for them regardless of which gene was
perturbed. The all-confidence tier does not measure regulatory recovery and we do not read it as
such.

Two conclusions follow. The absence of detectable transcription-factor specificity in the primary
CRISPRi panel is substantially a property of that assay: with a median of one measurable target per
factor, neither these features nor differential expression on the same cells nor established
network inference can demonstrate specificity there. On perturbation data that measure the
transcriptome the features do produce enrichment for curated regulons, at a rate above the
non-regulatory controls under the high-confidence reference and below them under the permissive
one. The high-confidence comparison is the interpretable one and it is under-powered: three factors
against one control gene, on a panel of thirteen. What can be stated is that these data do not
reproduce the flat null obtained on the primary panel, and that establishing whether a small
factor-specific signal exists would require a substantially larger panel than any single published
screen currently supports.


% Figure blocks. Each guards its own graphic so the manuscript still compiles while a
% figure is being regenerated; check_refs.py fails the build if any guard is still active
% at submission time.
\newcommand{\reviewfig}[4]{%
  \begin{figure}[htbp]\centering
  \IfFileExists{figures/#1}{\includegraphics[width=\textwidth]{figures/#1}}%
    {\fbox{\parbox[c][3cm][c]{0.9\textwidth}{\centering FIGURE PENDING: \texttt{\detokenize{#1}}}}}
  \caption{#2}\label{#3}
  \end{figure}#4}

\reviewfig{gb_fig1_atlas_overview.pdf}{\textbf{Sparse autoencoder feature atlas of Geneformer
V2-316M.} TopK sparse autoencoders trained on all 18 layers yield 82,525 alive features.
Reconstruction quality declines with depth while dead features increase.}{fig:overview}{}

\reviewfig{gb_fig4_ontology_heatmap.pdf}{\textbf{Biological annotation follows a U-shaped profile
across layers.} Early layers encode molecular programs, middle layers develop more abstract
representations, later layers re-specialize before becoming prediction-focused. Enrichment uses
Gene Ontology, KEGG (Kanehisa Laboratories), Reactome, STRING and TRRUST gene
sets.}{fig:ushape}{}

\reviewfig{gb_fig5_coactivation.pdf}{\textbf{Co-activation module structure across layers.}
(A)~Modules per layer. (B)~Fraction of alive features covered by a module. (C)~Co-activation edge
density.}{fig:modules}{}

\reviewfig{fig_highways.pdf}{\textbf{Cross-layer information flow is continuous.} (a)~Highway
fraction for the three long-range layer pairs after permutation-null correction; the
99th-percentile null threshold sits below the applied threshold for every pair, so the correction
does not shift the fractions. (b)~The same quantity for all 17 consecutive layer pairs, with no
sharp drop at any transition.}{fig:highways}{}

\reviewfig{fig_svd_capacity.pdf}{\textbf{The dictionary re-parameterizes the principal subspace
rather than reaching past it.} (a)~Cumulative projection $\rho_k$ of each decoder direction onto
the rank-$k$ principal subspace, median over features, against the random-direction expectation
$k/d$; the shaded band is the interquartile range at layer~17 and the dotted line marks the
50-axis basis. (b)~Truncated-SVD rank required to match the SAE's own variance explained, per
layer. (c)~Significant enrichment terms per direction when the leading 50 principal axes are
annotated with the protocol used for SAE features. Enrichment uses Gene Ontology, KEGG
(Kanehisa Laboratories) and Reactome gene sets.}{fig:svd}{}

\reviewfig{fig_concepts.pdf}{\textbf{Dictionary atoms and the distinct programs they express.}
(a)~Atoms and distinct gene-level programs per layer. (b)~Aggregate across all 18 dictionaries.
(c)~Mean absolute coherence between unit-normalized decoder atoms against the Welch bound, the
smallest achievable maximum coherence for that many unit vectors in 1,152
dimensions.}{fig:concepts}{}

\reviewfig{fig_causal.pdf}{\textbf{Causal feature ablation with an independent evaluation set.}
Distribution of specificity ratios across features, for three feature-selection rules and three
evaluation gene sets: the feature's top-20 genes intersected with its annotated term; the annotated
term's genes excluding that top-20; and a size-matched random gene set stratified on expression rank
and detection frequency, drawn per feature. The dotted line marks a ratio of one.}{fig:causal}{}

\reviewfig{fig_ceiling.pdf}{\textbf{What these data can support.} (a)~Fraction of transcription
factors whose predicted target set is enriched for its curated targets, at matched prediction-set
size and under one hypergeometric framework, for every method. (b)~How many perturbed
transcription factors have enough curated targets among the measured genes to be evaluable at
all.}{fig:ceiling}{}

\reviewfig{fig_power_replication.pdf}{\textbf{Unit of replication, sample size, and replication
across cell lines.} (a)~Responding features per target under a cell-level test against a
position-pooled test on the same cells; the dashed line is equality. (b)~Detection of the
knocked-down gene itself, and the number of responding features, as a function of perturbed cells
per target, restricted to the 14 factors with at least 100 cells so that sample size is not
confounded with which factors are available. (c)~Fraction of evaluable transcription factors with
target enrichment in each cell line, with perturbed and control cells drawn from the same
line.}{fig:power}{}

\section{Discussion}
We have presented a systematic application of sparse autoencoders to single-cell foundation models,
producing feature atlases for two architecturally distinct models: 82,525 features across 18 layers
of Geneformer V2-316M and 24,527 features across 12 layers of scGPT whole-human. Both atlases
reveal models that have organized biological knowledge into modular, interconnected internal
representations. Against that organization stands a specific negative result: the evidence that they carry
transcription-factor-to-target regulatory logic is bounded above by what the available perturbation
assays can measure, and on the assay this atlas was built from that bound is close to zero.

\paragraph{Superposition is real but is not invisibility.}
Comparing thousands of dictionary atoms with a handful of principal axes cannot establish that a
model's structure is inaccessible to linear decomposition, because the two sides of that comparison
differ in capacity by two orders of magnitude. Under a matched comparison the picture is sharper and
more interesting. SAE directions are not principal axes---even against the complete
1,152-dimensional basis, under 1\% of features align with any single axis above a cosine of 0.5---but
neither are they hidden from the principal subspace: the median direction places half of its norm
in the leading 150--250 components, well ahead of the random-direction reference, and truncated SVD
needs rank 250--403 to match the SAE's own reconstruction quality. Nor is biological annotation the
exclusive property of non-aligned directions: annotated with an identical protocol, the leading 50
principal axes yield at least as many enriched terms per direction as randomly drawn SAE features.
What the dictionary provides is not access to otherwise unreachable information but a
\emph{different parameterization} of it---many sparse, one-sided, individually interpretable
directions in place of a few dense, signed, polysemantic ones. That is a real and useful property,
and it is the one worth claiming.

\paragraph{What a feature count does and does not mean.}
Dictionary atoms trained at different layers occupy different vector spaces and cannot be summed
into a count of concepts coexisting in one representation. Within a single layer the dictionary is
4$\times$ overcomplete and resolves roughly 1,250 distinct gene-level programs from 4,608 atoms; in
aggregate the 18 dictionaries hold 82,525 atoms that resolve 18,711 distinct programs, with 97\% of
multi-atom programs recurring in two or more layers. The interesting quantity is not a compression
ratio but this pattern of re-derivation: each layer rebuilds a largely fresh dictionary, yet the
gene-level programs those dictionaries express recur throughout the stack.

\paragraph{Hierarchical biological abstraction across layers.}
The U-shaped annotation profile, the shift in module themes from molecular machinery to integrative
programs, and the near-complete representational turnover between adjacent layers together describe
hierarchical biological abstraction. This parallels findings in large language models, where early
layers handle token-level processing and later layers more abstract semantic computation, but here
the tokens are genes and the semantics are biological programs.

\paragraph{Feature-level structure versus component-level nulls.}
Ablating a single feature produces a large, reproducible and gene-specific change in the model's
output where ablation of whole attention heads or MLP layers produces null behavioural
effects~\cite{kendiukhov2025systematic}. Computation is therefore organized around directions in
the residual stream rather than around architectural components, which is why component-level
interventions missed it. The scope of that effect is narrow, however: it falls on the genes the
feature fires on and not on the remainder of the biological process the feature is annotated with,
which is indistinguishable from a size- and activation-matched random gene set. A feature is a
handle on its own support, not on the pathway its label names---a distinction that matters for
anyone intending to use SAE features as intervention targets.

\paragraph{The co-expression--regulation boundary, and where it actually lies.}
Our results extend to the residual stream a boundary that has now been located by several
independent probes of the same models. Attention-derived edge scores add no incremental predictive
value over gene-level features for perturbation-target prediction~\cite{kendiukhov2025systematic}.
Causal circuit tracing recovers a dense, biologically coherent computational graph whose edges do
not correspond to directed regulatory relationships~\cite{kendiukhov2026circuittracing}, and
exhaustive mapping of that graph finds massive redundancy rather than identifiable regulatory
pathways~\cite{kendiukhov2026exhaustive}. Direct causal intervention on gene representations in
scGPT does not recover regulatory targets~\cite{kendiukhov2026causalintervention}, and adversarial
validation of in-silico perturbation profiles shows the same for predicted perturbation
responses~\cite{kendiukhov2026insilico}. Analyses of the representation geometry reach the same
place from a different direction: the structure these models learn is organized by cell state and
pathway membership~\cite{kendiukhov2026topology,kendiukhov2026celltypeprograms}, and although
residual-stream geometry does carry an incremental gene-regulatory signal, that signal is fragile
and does not survive stricter cross-validation~\cite{kendiukhov2026grsgeometry}. Sparse
dictionaries are one more probe, applied to the one representation the others did not decompose,
and they find what the others found.

No finite set of probes can establish that regulatory information is absent from a model, because
the proposition that it resides in some representation not yet examined is not refutable by
examination. What can be done is to state which representations have been examined, with which
instruments, and what was found. Attention weights, individual components, the causal graph
connecting them, gene-embedding geometry, in-silico perturbation responses, and now a sparse
decomposition of the residual stream have each been examined; none yields transcription
factor-to-target structure above what the underlying measurements support on the assays where each
was applied. That is the sense in
which we take the finding to be a property of these models rather than of any one instrument.

The important qualification is what that boundary can be attributed to, and here the evidence is
two-sided in a way that the primary dataset alone could not reveal.

A specificity rate is uninterpretable without knowing how much of the regulatory relationship is
recoverable from the same measurements. On the CRISPRi resource the atlas was built from, that
ceiling is very low: the assay measures 6,546 genes and the median perturbed transcription factor
has a single curated target among them. Differential expression computed directly on the knockdown
recovers curated targets for a minority of factors, established network inference recovers no more,
and the sparse features recover targets for 1 of 66 factors across four cell lines. The correct reading of that panel
is not that the models lack regulatory structure but that the measurement cannot decide the
question.

On perturbation data that measure the transcriptome the flat null does not reproduce. In a
CRISPR-activation screen with 33,694 genes measured and strong induction, a factor's predicted
target set is enriched for its own high-confidence regulon in 3 of 13 cases, against 1 of 9
perturbations of genes absent from every reference network. That is the direction specificity
predicts, and the panel is too small to establish it ($p = 0.62$). Against low-confidence regulons
of several hundred genes the comparison reverses outright---8 of 9 controls enrich---so that tier
measures breadth of the predicted set rather than regulatory recovery.

What this study can support is therefore bounded on both sides. It is not the absence of regulatory
information: the assay the atlas was built on cannot detect such information if it is present, and
on data that can, the null does not hold. Neither is it the presence of a factor-specific signal:
the one comparison capable of demonstrating that rests on thirteen factors in a single screen and
does not reach significance. Settling the question needs a perturbation panel substantially larger
than any single published screen provides. What we can state without qualification is that an
interpretability result reported on a 6,546-gene assay is measuring the assay as much as the
model.

\paragraph{Where the unit of analysis matters, and by how much.}
A perturbation is applied to a cell, but a per-position statistic treats every gene position in
that cell as a separate observation. Whether this matters depends on how strongly feature
activation is correlated within a cell, which is an empirical question we can answer rather than
assume. The intraclass correlation of SAE feature activation across positions of the same cell is
$2.98 \times 10^{-4}$, which for the observed 2,020 positions per cell gives a design
effect of 1.61, that is an effective sample size of 62\% of the nominal count. The effective sample size is therefore a modest fraction of the
nominal position count rather than smaller by three orders of magnitude: sparse feature activations
at different gene positions of one cell are close to independent, because the sparsity pattern
dominates the within-cell variance.

The correction is nonetheless the right one to make, and we make it: significance is assessed with
the cell as the unit of replication throughout, holding the effect-size definition fixed at the
position-level scale so that the cell-level and position-level results differ only in what is
counted as an independent observation. We recommend the cell-level formulation as the default for
any interpretability analysis that tests per-position statistics against a per-cell intervention,
while noting that its practical effect here is bounded and measurable.

\paragraph{Implications for foundation model training.}
Current scFM training objectives may inherently bias representations toward co-expression. Learning
causal regulatory relationships would require training signals that distinguish cause from
correlation, such as perturbation-prediction objectives incorporated during pre-training. Equally,
evaluating whether a model has learned such relationships requires assays that measure the target
genes: a genome-scale perturbation atlas with a restricted expression panel bounds what any
interpretability method can demonstrate, however good the method is.

\paragraph{SAEs as a general interpretability framework for biological models.}
Independent of the regulatory question, this work establishes SAEs as a productive framework for
scFM interpretability. Applying an identical pipeline to two architecturally distinct models, with
qualitatively consistent results, demonstrates that these tools transfer to any transformer-based
biological model.

\paragraph{What would settle it.}
Regulatory structure in these representations should be assessed on perturbation atlases that
measure the whole transcriptome, against high-confidence regulons, with perturbations of
non-regulatory genes as a matched control, and with the recoverable ceiling computed on the same
cells. We apply that design here, and its limitation is scale rather than construction: thirteen
evaluable factors in one screen cannot separate a modest specific signal from none. A panel several
times larger, or several screens combined under a common evaluation, would. Applied instead to a
6,546-gene assay, any interpretability method returns a null whose cause cannot be identified, and
we would encourage reporting the evaluability of the assay alongside any claim about what a model
encodes.

\paragraph{Limitations.}
Several bounds apply to what we report. The regulatory analysis is limited by the gene panel of the
perturbation assay, as quantified above; a perturbation resource measuring the full transcriptome
would raise the ceiling against which any method is judged. TRRUST and DoRothEA are themselves
incomplete and literature-biased; a ChIP-seq-derived catalog would give an orthogonal view we have
not pursued. Our SAEs use one point in a wide architecture space, and although a targeted
expansion-by-sparsity ablation and a multi-seed stability analysis both leave the conclusions
intact, the space is not exhausted. Causal patching tests single-feature ablation; combinatorial
effects remain unexplored. Our SAEs are trained on unperturbed activations, which may bias the
learned dictionary toward cell-state and co-expression programs; training on pooled perturbed and
control activations with perturbation identity as an auxiliary signal is the most direct way to test
this and is beyond the scope of the present analysis. Finally, donor and assay identity are not
recorded in the cached extraction metadata for the Tabula Sapiens arm, so the batch analysis rests
on tissue as a coarse proxy.

\section{Methods}
\subsection*{Models and data}

\paragraph{Geneformer.} We used Geneformer V2-316M~\cite{theodoris2023transfer} (18 transformer
layers, 1,152 hidden dimensions, 18 attention heads per layer) from HuggingFace
(\texttt{ctheodoris/\allowbreak Geneformer}, subfolder \texttt{Geneformer-V2-316M}). Perturbation
data come from the Replogle genome-scale CRISPRi resource~\cite{replogle2022mapping} as a single
processed matrix spanning four cell lines (K562, RPE1, Jurkat, HepG2; 643,413 cells, 6,546 measured
genes, 2,023 perturbation targets plus non-targeting controls).

\paragraph{Control cohort.} The atlas was built from 2,000 non-targeting control cells sampled at
random (seed 42) from the 39,165 non-targeting cells in the resource. Because the control condition
is defined by perturbation label, this cohort spans all four lines: 591 Jurkat, 581 K562, 567 RPE1
and 261 HepG2 cells (mean 2,028 genes per cell, 4,056,351 total token positions). The atlas
dictionaries are therefore learned from a multi-cell-line control population, and we refer to them
throughout as the multi-line control atlas rather than attributing them to any single line. Cohort
composition for every analysis is recorded in \suptab{cohorts}. For the multi-tissue
comparison we additionally used 3,000 cells from Tabula Sapiens~\cite{tabula2022tabula} (1,000
immune cells across 43 cell types, 1,000 kidney cells, 1,000 lung cells), yielding 5,623,164 token
positions.

\paragraph{Cell-line assignment in downstream analyses.} Analyses that compare perturbed against
control cells draw both from the same line, so that a difference between the two cannot reflect a
difference in cell-line composition. Causal patching is likewise run within a single stated line.
Where an analysis uses the atlas dictionary on cells of one line, this is a deliberate transfer
setting and is labelled as such.

\paragraph{scGPT.} We used scGPT whole-human~\cite{cui2024scgpt} (12 transformer layers, 512 hidden
dimensions, 8 attention heads per layer), trained on approximately 33 million single-cell
transcriptomic profiles. scGPT ingests each cell as a pair of parallel token sequences:
gene-identity tokens drawn from a vocabulary of approximately 60,700 genes, and \emph{binned}
expression-value tokens (log1p-normalized counts discretized into 51 bins by default). Genes are
sorted by expression level (descending) and padded or truncated to a maximum sequence length of
1,200. Pre-training uses an attention-masking scheme (``gene prompt'' and ``value prompt'' modes)
that enables generative prediction of expression values conditional on gene identity, which we
refer to throughout as a generative or causal-attention objective to distinguish it from
Geneformer's bidirectional masked-token objective. Activations were extracted from the same 3,000
Tabula Sapiens cells, yielding 3,561,832 token positions per layer.

\paragraph{Input handling for each model.} The two models differ at the first layer, and this
difference propagates through all downstream analyses. For Geneformer the input is deterministic
given a cell: the per-gene expression vector is sorted, each gene receives a rank token, and no
continuous magnitude is passed to the model. For scGPT both the gene-identity token and its bin
index are embedded and summed before the transformer stack, so the model has direct access to a
coarse magnitude signal. Residual-stream activations therefore capture, at every layer, a
representation whose only expression signal is gene \emph{rank} for Geneformer, and one that
includes both rank-ordered gene identity and a binned magnitude for scGPT.

\subsection*{Activation extraction}

For both models we performed forward passes with PyTorch hooks registered at each transformer
layer's output (after the residual connection), collecting hidden states at every gene position.
Activations were stored as memory-mapped NumPy arrays (float32). Extraction used Apple Silicon MPS
acceleration. For Geneformer ($d{=}1{,}152$) this is 336.4~GB for the 18 control-cohort layers (18.7~GB per
layer) and 103.6~GB for four Tabula Sapiens layers; for scGPT ($d{=}512$), approximately 82~GB for
12 layers.

For the perturbation analyses, where the quantity of interest is a per-cell summary rather than the
raw residual stream, we instead hooked the target layer directly, encoded the captured hidden state
through the layer's SAE on the accelerator, and reduced to per-cell statistics (mean activation,
firing rate, and mean square of each feature over that cell's gene positions) before transfer.
Cells were batched by sequence length. This keeps the cell as the retained unit of observation and
makes the per-cell feature matrices small enough to resample freely.

\subsection*{TopK sparse autoencoder architecture and training}

We used TopK SAEs~\cite{makhzani2013ksparse} parameterized by input dimension $d$. The encoder maps
$\mathbf{x} \in \mathbb{R}^{d}$ (centered by subtracting the training-set mean $\boldsymbol{\mu}$)
through a linear projection to a pre-activation vector in $\mathbb{R}^{4d}$, followed by TopK
sparsification retaining only the $k{=}32$ largest activations, to which a ReLU is applied. The
decoder reconstructs via $\hat{\mathbf{x}} = W_\text{dec}\mathbf{h}_\text{sparse} +
\mathbf{b}_\text{dec}$, where $W_\text{dec}$ columns are re-normalized to unit norm after each
gradient step. Training minimized reconstruction MSE; sparsity is enforced structurally by the TopK
operation rather than by an $L_1$ penalty.

For Geneformer ($d{=}1{,}152$): 4,608 features per layer, trained on a 1M-position subsample per
layer. For scGPT ($d{=}512$): 2,048 features per layer, trained on all 3,561,832 positions per
layer. Both used Adam, learning rate $3 \times 10^{-4}$, batch size 4,096, and 5 epochs. The atlas
SAEs were trained with random seed 42, which also fixes the training subsample; the sensitivity of
every reported quantity to this choice is characterized in the seed-stability analysis below.

\subsection*{Feature analysis and ontology annotation}

For each alive feature (activation frequency $> 0$ on 100K held-out positions) we identified the
top 20 genes by mean activation magnitude. Each feature was tested for enrichment against Gene
Ontology Biological Process~\cite{ashburner2000go}, KEGG pathways~\cite{kanehisa2000kegg}, Reactome
pathways~\cite{jassal2020reactome}, STRING protein--protein interactions~\cite{szklarczyk2023string},
and TRRUST TF--target relationships~\cite{han2018trrust}. TF--target priors from
DoRothEA~\cite{garcia-alonso2019dorothea} were used at the A+B+C confidence subset and at the union
of all confidence levels. Enrichment was assessed with a one-sided Fisher's exact test and
Benjamini--Hochberg FDR at $\alpha = 0.05$, with term sizes restricted to 5--500 genes.

\subsection*{Comparison against the principal subspace}

At each layer we computed the exact eigendecomposition of the activation covariance from a
500,000-position sample (40 evenly spaced contiguous blocks), giving all $d$ principal directions
and their eigenvalues. Three comparisons follow.

\emph{Reconstruction-matched capacity.} The truncated-SVD rank $k^{*}$ at which cumulative
eigenvalue mass equals the SAE's variance explained on held-out positions.

\emph{Cumulative projection.} For each unit-normalized decoder direction $\mathbf{d}_f$ we computed
$\rho_k = \lVert P_k \mathbf{d}_f \rVert^2 = \sum_{i \le k} (\mathbf{d}_f \cdot \mathbf{v}_i)^2$,
the fraction of the direction's norm lying in the rank-$k$ principal subspace, for all $k$ from 1
to $d$. The reference for these curves is a random unit direction, for which
$\mathbb{E}[\rho_k] = k/d$.

\emph{Single-axis alignment.} The fraction of features whose absolute cosine similarity with at
least one of the leading $k$ axes exceeds $\tau$, reported jointly over
$k \in \{1, \dots, d\}$ and $\tau \in \{0.3, 0.5, 0.7, 0.9\}$ so that the figure does not depend on
either choice.

\emph{Annotation yield per direction.} The leading 50 principal axes were annotated with the
identical top-20-gene enrichment protocol used for SAE features. Because a principal axis is signed
whereas a TopK feature is one-sided, each axis was annotated in both polarities (genes with the
most positive and most negative mean projection), and results are reported per polarity and for
their union.

\subsection*{Counting distinct concepts across layers}

Dictionary atoms from different layers occupy different vector spaces and cannot be pooled as
directions. They can be compared in the shared gene space through their top-20 gene signatures. We
built a similarity graph over all 82,525 Geneformer features, placing an edge between two features
whose top-20 gene sets share at least 5 genes and have Jaccard similarity at least 0.14, and
partitioned it with the Leiden algorithm~\cite{traag2019louvain} at resolutions
$\gamma \in \{0.5, 1.0, 2.0\}$. The number of communities is the number of distinct gene-level
programs. We additionally report, per layer, the number of atoms, the number of distinct programs,
the mean absolute decoder coherence, and the Welch bound
$\sqrt{(n-d)/(d(n-1))}$ for $n$ unit vectors in $\mathbb{R}^{d}$, which is the smallest achievable
maximum coherence and therefore the correct reference for how many near-orthogonal directions a
$d$-dimensional space can host.

\subsection*{Co-activation graph and module detection}

For each layer, pointwise mutual information (PMI~\cite{manning1999foundations}) was computed
between all pairs of alive features across the layer's token positions; feature $i$ is active at
position $j$ if it is among the top-$k$ activations there, and
$\text{PMI}(i,j) = \log_2 P(i,j) / [P(i)P(j)]$. Edges were retained where two features co-fired at
a minimum of 50 positions with PMI at least 2.0 bits. Community detection used the Leiden algorithm~\cite{traag2019louvain} at
resolution $\gamma = 1.0$, the \texttt{leidenalg} default. Because a single resolution is a choice
rather than a ground truth, we swept $\gamma \in \{0.5, 0.75, 1.0, 1.5, 2.0\}$ at layers 0 and 11
and report the resulting partition sizes and adjusted Rand indices; module counts in the main text
are the $\gamma{=}1.0$ cut of a hierarchical structure whose qualitative module identities are
recognizable throughout the range.

\subsection*{Causal feature patching}

Single-feature ablation was performed with PyTorch forward hooks. At the target layer's output the hidden
state was encoded through the SAE and the contribution of one feature, its activation times its
decoder atom, was subtracted from the residual stream; the forward pass then continued to the
output. Because the intervention is subtractive rather than a substitution of the reconstruction,
no reconstruction error is introduced. The readout is the logit of each position's own gene token. Disruption is the
mean change in output logits at a set of gene positions.

Three properties of the design keep the evaluation independent of the annotation. First, features
are drawn from three arms: the most richly annotated features, a uniformly random sample of
annotated features, and a uniformly random sample of all alive features. Second, for each feature we
partition gene positions into those whose gene is in both the feature's top-20 activating genes and
its annotated term, those in the annotated term but \emph{not} in the top-20 (the independent
evaluation set), and all others. Third, for every feature we draw a random gene set matched in size to the independent evaluation
set and stratified on gene expression rank (ten bins) and detection frequency across the patched
cells (three bins), drawn from outside both the annotated term and the feature's top-20 genes,
giving each feature its own null.
Specificity is the ratio of mean absolute disruption on a gene set to that on all other positions,
and the held-out ratio is reported against the matched-random ratio for the same feature.

\subsection*{Perturbation response mapping}

For each perturbation target we sampled up to 200 CRISPRi cells and compared them against 600
non-targeting control cells of the same line, encoding every cell through the layer-11 SAE as
described above. A feature responds if its per-cell mean activation differs between perturbed and
control cells (Mann--Whitney $U$ over cells, Benjamini--Hochberg across features, FDR $< 0.05$) and
the shift exceeds half a control standard deviation, the latter measured on the position-level
scale so that the cell-level and position-level analyses differ only in the inference. The cell is the unit of replication throughout: gene positions
within a cell are not independent observations, and we report the intraclass correlation of feature
activation within cells together with the number of features that a position-pooled test would call
significant on exactly the same cells.

TF specificity is evaluated in the same framework as every other method (below): the union of the
top-20 gene lists of responding features, ranked by effect size, forms the predicted target set,
which is tested for enrichment of the TF's curated targets by hypergeometric test with
Benjamini--Hochberg correction across the panel. A label-permutation null, in which perturbed and
control labels are shuffled across the same cells, calibrates the procedure.

\subsection*{Recoverability benchmark}

A recovery rate is interpretable only against the amount of information the measurements contain.
We therefore evaluated several methods on one footing: each emits a ranked list of putative targets
for every TF in the panel, the top 100 of which are tested for enrichment of that TF's curated
targets by hypergeometric test with Benjamini--Hochberg correction across the panel.

\emph{Panel.} TFs perturbed in the cell line with at least 50 cells and at least 5 curated targets
present among the 6,546 measured genes. The second criterion is reported explicitly, because a TF
whose curated targets are not measured cannot show target-specific behaviour under any method.

\emph{Universe.} The hypergeometric universe is the set of measured genes. Using a nominal
20,000-gene universe is anticonservative here, since every method can only draw predictions from
measured genes; we report both, and the random-gene-set control is calibrated only under the
measured-gene universe.

\emph{Methods.} (i)~Differential expression of the knockdown itself (Mann--Whitney over cells,
perturbed vs non-targeting, Benjamini--Hochberg across genes), which is perturbation-aware and
constitutes the empirical ceiling attainable on these data. (ii)~GENIE3~\cite{huynhthu2010genie3}
and (iii)~GRNBoost2~\cite{moerman2019grnboost2}, in their published formulation: every target gene
is regressed on the candidate regulators with a tree ensemble and regulators are ranked by variable
importance (random forest with \texttt{max\_features}~=~$\sqrt{p}$; and stochastic gradient boosting
with learning rate 0.01, \texttt{max\_features}~=~0.1, subsample 0.9 and early stopping), fitted on
the same control cells. (iv)~Pearson and (v)~Spearman co-expression with the TF on the same control
cells. (vi)~Cosine similarity between mean contextual gene embeddings taken from the Geneformer
residual stream, a model-internal baseline that uses no SAE. (vii)~Size-matched random gene sets.
(viii)~The SAE feature response described above.

\subsection*{Replication across cell lines}

The perturbation analysis was repeated in RPE1, Jurkat and HepG2 cells from the same CRISPRi
resource, which are independent screens in independent cell backgrounds. These lines were run at
100 cells per target against 400 non-targeting cells over 16 factors and 6 non-regulatory controls,
against 200, 600, 20 and 10 respectively for the primary line. In
every line, perturbed cells are compared against non-targeting cells of that same line. Each line
was analysed with two dictionaries: the multi-line control atlas dictionary, and the dictionary
additionally trained on Tabula Sapiens primary tissue, which asks whether widening the training
distribution beyond immortalized lines changes what is recovered.

\subsection*{Independent external perturbation datasets}

The measurement was additionally repeated on two perturbation datasets generated by other groups.
ECCITE-seq of THP-1 monocytes~\cite{papalexi2021eccite} measures 18,649 genes; its perturbation
labels are guide-level and were reduced to the targeted gene. Perturb-seq of
K562~\cite{norman2019exploring} measures 33,694 genes and is a CRISPR-activation screen, so it
tests the same question under gene induction rather than repression; combinatorial perturbations
were excluded. Counts were CP10K-normalized and log1p-transformed as for the primary resource, and
genes were mapped to the model vocabulary by symbol, which covers 13,601 of 18,649 and 18,121 of
33,694 genes respectively.

Each dataset was analysed with the atlas dictionary at layer~11 under the pipeline described above:
100 cells per target against 400 non-targeting cells, cell-level Mann--Whitney with
Benjamini--Hochberg across features, and the same hypergeometric enrichment of a 100-gene predicted
target set. Panels comprise the perturbed transcription factors with at least 50 cells and at least
five curated targets among the measured genes.

For the specificity control, ten perturbations of genes absent from every reference network were
drawn at random from the same dataset and run through the identical pipeline. Their predicted sets
were then tested against the \emph{same} regulons as the transcription factors, giving one
enrichment test per control-gene-by-regulon pair, with Benjamini--Hochberg applied within each arm.
Because a control gene contributes many tests, the comparison between arms is made at the level of
independent perturbations---how many control genes enrich for at least one regulon against how many
transcription factors enrich for their own---by Fisher's exact test.

\subsection*{Run-to-run stability}

The atlas SAE was retrained at layers 0 and 11 under five independent seeds with the
training subsample held fixed, isolating initialization and batch-order stochasticity, plus one
additional run per layer with a different 1M-position training subsample. Dictionaries were compared
by decoder cosine similarity, reporting both the unconstrained best match for each atom and a greedy
one-to-one assignment. Variance explained, dead-feature count, decoder coherence, annotation rate,
and module count were recomputed per run.

\subsection*{Matched cross-model comparison}

Because the two atlases were built from different inputs, architecture and input distribution are
confounded in any direct comparison of their statistics. We therefore trained Geneformer SAEs on the
same Tabula Sapiens cells used for the scGPT atlas, with identical hyperparameters, and compared the
two models on matched data at matched relative depth (layer index divided by depth minus one). For
the Geneformer arm we additionally separate two effects by reporting the variance explained when the
control-cohort SAE is applied to Tabula Sapiens activations against that of the SAE trained on those
activations directly.

\subsection*{Interactive web atlases}

Both feature atlases were deployed as single-page web applications built with React 18, TypeScript,
Vite 6, Tailwind CSS, and Plotly.js. All feature data were preprocessed into compact JSON files
served statically via GitHub Pages.

\subsection*{Dimensionality reduction and visualization}

We projected decoder weight vectors to two dimensions using UMAP~\cite{mcinnes2018umap} with cosine
distance, obtaining structureless point clouds as expected from quasi-orthogonal directions. We
therefore visualized the co-activation graph directly with a force-directed layout
(Fruchterman--Reingold via NetworkX); the corresponding supplementary figures are force-directed
layouts of that graph and are not UMAP or t-SNE projections. For independent validation,
annotation-based feature vectors (TF-IDF weighted ontology terms) were projected with UMAP
($n_\text{neighbors}=15$, $\min_\text{dist}=0.1$, cosine metric) and t-SNE (perplexity 20, cosine
metric).

\subsection*{Computational environment}

All experiments were run on a single Apple Silicon workstation (Apple M2 Pro, 10 cores, 32~GB
unified memory). Acceleration used PyTorch's MPS backend, with CPU fallback for operations without
an MPS implementation (Leiden community detection, ontology Fisher tests, tree-ensemble regression).
Software stack: Python 3.11; PyTorch 2.x with MPS; NumPy, SciPy, scikit-learn, pandas, scanpy,
anndata, leidenalg, python-igraph, networkx, umap-learn, statsmodels, h5py. Exact package versions
are pinned in \texttt{requirements.txt} in the code repository. Wall-clock times for the longest
steps are logged alongside the compute environment: Geneformer control-cohort activation extraction
$\sim$8~h, SAE training $\sim$3~min per layer, per-layer PMI and Leiden $\sim$20--60~min,
per-cell perturbation encoding $\sim$0.9 cells~s$^{-1}$.

\section*{Data availability}
All datasets analysed in this study are publicly available.

\begin{itemize}\itemsep2pt
\item \textbf{Genome-scale CRISPRi Perturb-seq} (K562, RPE1, Jurkat, HepG2)~\cite{replogle2022mapping}:
Figshare+, \url{https://doi.org/10.25452/figshare.plus.20029387.v1}.
\item \textbf{ECCITE-seq of THP-1 monocytes}~\cite{papalexi2021eccite}: available through the
scPerturb collection, \url{https://doi.org/10.5281/zenodo.13350497}.
\item \textbf{Perturb-seq of K562}~\cite{norman2019exploring}: available through the scPerturb
collection, \url{https://doi.org/10.5281/zenodo.13350497}.
\item \textbf{Tabula Sapiens}~\cite{tabula2022tabula}:
\url{https://tabula-sapiens-portal.ds.czbiohub.org/} and CZ CELLxGENE.
\item \textbf{Geneformer V2-316M}~\cite{theodoris2023transfer}:
\url{https://huggingface.co/ctheodoris/Geneformer} (subfolder \texttt{Geneformer-V2-316M}).
\item \textbf{scGPT whole-human}~\cite{cui2024scgpt}: \url{https://github.com/bowang-lab/scGPT}.
\item \textbf{TRRUST v2}~\cite{han2018trrust}: \url{https://www.grnpedia.org/trrust/}.
\item \textbf{DoRothEA}~\cite{garcia-alonso2019dorothea}: \url{https://saezlab.github.io/dorothea/}.
\item \textbf{Gene Ontology}~\cite{ashburner2000go}: \url{http://geneontology.org/}.
\item \textbf{KEGG}~\cite{kanehisa2000kegg}: \url{https://www.kegg.jp/}. KEGG pathway gene sets are
used under permission from Kanehisa Laboratories.
\item \textbf{Reactome}~\cite{jassal2020reactome}: \url{https://reactome.org/}.
\item \textbf{STRING}~\cite{szklarczyk2023string}: \url{https://string-db.org/}.
\end{itemize}

\noindent
The two feature atlases produced by this work are released as interactive web platforms:
Geneformer (\url{https://biodyn-ai.github.io/geneformer-atlas/}) and scGPT
(\url{https://biodyn-ai.github.io/scgpt-atlas/}).

\section*{Code availability}
All analysis code, trained sparse autoencoder checkpoints, feature catalogs, and the result files
from which every figure and table in this article is generated are archived at Zenodo,
\textbf{[Zenodo DOI to be inserted at acceptance]}, and are also available at
\url{https://github.com/Biodyn-AI/bio-sae}. The archive includes the figure- and
table-generation scripts, so each reported number can be traced to the result file that produced
it, together with a reference-consistency check that verifies the correspondence between the main
text and the Supplementary Information. Source for the interactive atlases is at
\url{https://github.com/Biodyn-AI/geneformer-atlas} and
\url{https://github.com/Biodyn-AI/scgpt-atlas}.

\section*{Acknowledgements}
Computations were performed on Apple Silicon hardware with MPS acceleration.

\section*{Author contributions}
I.K. is the sole author. I.K. conceptualized the study, developed the methodology, wrote all
software, performed all computational experiments, built the interactive feature atlases,
analysed and interpreted the results, and wrote the manuscript.

\section*{Competing interests}
The author declares no competing interests.

\bibliography{references}

\end{document}
